% Generated draft: requires mathematical and source review.
\documentclass[11pt]{article}
\usepackage[margin=25mm]{geometry}
\usepackage{fontspec}
\setmainfont{Latin Modern Roman}
\usepackage{amsmath,amssymb,mathtools,unicode-math}
\setmathfont{Latin Modern Math}
\usepackage{enumitem,needspace,xurl,hyperref}
\setlength{\parindent}{0pt}
\setlength{\parskip}{5pt}
\setlength{\emergencystretch}{3em}
\newcommand{\N}{\mathbb{N}}
\newcommand{\Z}{\mathbb{Z}}
\newcommand{\Q}{\mathbb{Q}}
\newcommand{\R}{\mathbb{R}}
\newcommand{\C}{\mathbb{C}}
\newcommand{\F}{\mathbb{F}}
\newcommand{\A}{\mathbb{A}}
\newcommand{\PP}{\mathbb{P}}
\newcommand{\E}{\mathbb{E}}
\newcommand{\eps}{\varepsilon}
\newcommand{\dd}{\,\mathrm{d}}
\begin{document}
\section*{UnsolvedMath ID: 20001967}
\textbf{Source problem code:} catalogue entry
\subsection{Definitions and mathematical statement}
The main source is Conjecture \(6.16\) on the rank-\(n\), degree-\(d\), fixed-determinant trace-free Higgs moduli space \(M=\mathcal M_{\mathrm{Dol}}^d(\operatorname{SL}(n,\C))\) over a compact Riemann surface. Precise determinant, stability, coprimality and genus hypotheses are not given in the conjecture excerpt. With \(F\) denoting fixed components for the relevant Higgs-field scaling action, it writes the localized equivariant volume
\[
 \int_Me^\alpha=\sum_F\int_F\frac{e^{\alpha|_F}}{E(N_F)}.
\]
Here an equivariant Euler denominator and localized noncompact integration convention are intended; neither the equivariant class \(\alpha\) nor that convention is supplied. The conjecture asks for this and all equivariant intersection numbers as iterated residues of an expression \(BA(y_1,\ldots,y_n)\). Poles indexed by ordered partitions \(n=\lambda_1+\cdots+\lambda_k\) should reproduce fixed-locus contributions of type \((\lambda_1,\ldots,\lambda_k)\); all remaining poles should solve Bethe Ansatz equations, with
\[
 \int_Me^\alpha=-\sum_{\text{Bethe poles}}\operatorname{Res}(BA).
\]
The expression, equations, independent variables, residue order, orientations, Weyl factors and compactification are absent, so this is the source's schematic target, not a fully specified residue identity.

The aid reports a rigorous rank-two formula for a universal invariant sector, but no all-rank proof of this exact conjecture. Its all-rank conditional reduction requires an applicable compact several-variable global-residue theorem; a partition-pole/fixed-locus match with consistent residue order and Weyl normalization; exhaustion of finite poles by partition and Bethe poles; and vanishing total boundary residue. Under those checks the global zero-sum identity gives the displayed negative Bethe sum. The checks must also hold for each insertion to cover all intersection numbers. They are not proved for the unspecified \(BA\) here.

For collision stability the aid uses the following complete-intersection residue formulation. Let \(f_1(y,t),\ldots,f_r(y,t)\) and \(a(y,t)\) be holomorphic on a parameterized neighborhood of a relatively compact domain in \(\C^r\). Assume the common-zero scheme is finite, isolated inside the domain, with no zeros escaping or meeting its boundary and with the local complete-intersection residue setup preserved. For simple zeros at generic \(t\), the total residue is
\[
 \sum_{f(y,t)=0}\frac{a(y,t)}{\det(\partial f_i/\partial y_j)(y,t)}.
\]
Its limit when roots collide is the sum of the local Grothendieck residues of
\(a\,\dd y_1\wedge\cdots\wedge\dd y_r/(f_1\cdots f_r)\)
on the entire zero-dimensional scheme at the limiting parameter, including multiplicities. This defines a collision-stable total Bethe contribution under the stated no-escape and residue hypotheses. An individual Jacobian denominator is not used at a multiple root; positive-dimensional limits or boundary escape are outside this statement.

The canonical extraction also contains Section \(6.5\), ``Arithmetic approach.'' Its data and normalization issues are retained separately. For a compact Riemann surface \(\Sigma\) of genus \(g\), integers \(n\geq1,d\) with \(\gcd(n,d)=1\), and primitive central twist \(\xi_n\) corresponding to \(d\), it describes
\[
 \mathcal M_B^d(\operatorname{GL}(n,\C))=
 \{(A_i,B_i)\in\operatorname{GL}(n,\C)^{2g}:
       \prod_{i=1}^g A_i^{-1}B_i^{-1}A_iB_i=\xi_n I\}/\operatorname{GL}(n,\C)
\]
and Simpson's nonabelian-Hodge diffeomorphism with the rank-\(n\), degree-\(d\) Dolbeault moduli. Over \(\F_q\) the twist must exist, in particular with characteristic prime to \(n\) and a suitable \(n\)-th root of unity in the field. Quotient means a moduli/GIT construction; orbit, coarse-point and stack-weighted counting must be distinguished.

The source's Frobenius--Schur character expression is
\[
 S_{n,g}(q)=\sum_{\chi\in\operatorname{Irr}(\operatorname{GL}(n,\F_q))}
       \frac{|\operatorname{GL}(n,\F_q)|^{2g-2}}{\chi(1)^{2g-1}}
       \chi(\xi_n I).
\]
It defines the ordinary mixed-Hodge polynomial
\(\mathcal H_n(x,y,t)=\sum_{p,q,k}h^{p,q;k}x^py^qt^k\)
from Deligne's filtrations and labels as Theorem \(6.17\) the equality of \(\mathcal H_n(\sqrt q,\sqrt q,-1)\), the finite-field quotient count, and \(S_{n,g}(q)\). This equality needs a counting audit: already for \(n=1\), the coarse quotient count is \((q-1)^{2g}\), while the printed character sum is \((q-1)^{2g-1}\). The residual scalar-stabilizer factor separates coarse and weighted counts. Point-count/Hodge comparison also requires the appropriate arithmetic spread and support/normalization conventions; the printed equality is therefore recorded as a source theorem claim with an unresolved normalization issue, not silently corrected.

The same arithmetic passage introduces formal two-variable functions, distinct from \(\mathcal H_n(x,y,t)\):
\[
 V_n(q,t)=\frac{H_n(q,t)(qt^2)^{(1-g)n(n-1)}}{(qt^2-1)(q-1)},\qquad
 Z_n(q,t,T)=\exp\!\left(\sum_{r\geq1}
          \frac{V_n(q^r,-(-t)^r)T^r}{r}\right),
\]
\[
 H_\lambda^g(q,t)=\prod_{z\in d(\lambda)}
 \frac{(qt^2)^{(2-2g)\ell(z)}
       (1+q^{h(z)}t^{2\ell(z)+1})^{2g}}
      {(1-q^{h(z)}t^{2\ell(z)+2})(1-q^{h(z)}t^{2\ell(z)})},
\]
\[
 \prod_{n\geq1}Z_n(q,t,T^n)
       =\sum_{\lambda\in\mathcal P}H_\lambda^g(q,t)T^{|\lambda|}.
\]
Here \(\mathcal P\) is the set of partitions and \(d(\lambda)\) its Young diagram. These formulas retain the intended fraction layout of the damaged extraction. It calls \(a(z),\ell(z),h(z)\) arm, leg and hook lengths and prints \(h=a+\ell-1\). Under the usual convention excluding the box from arm and leg, the hook is instead \(a+\ell+1\); inclusive conventions and the diagram must be checked before using these formulas. The excerpt ends before stating its promised subsequent conjectures, so no further Hodge-polynomial identity is invented.
\subsection{Short English statement}
The Bethe-residue conjecture is supplied only schematically. The aid reduces it to explicit global-residue and localization checks and treats colliding Bethe roots by total scheme residues. The merged arithmetic passage is retained with its character-count normalization discrepancy, hook-length ambiguity and missing later conjectures.
\subsection{Sources}
\begin{itemize}
\item Local UnsolvedMath registry, permanent ID \(20001967\), source code \texttt{catalogue entry}; canonical collection \texttt{aim-geometry-notes.json}, notes[304].
\item AIM workshop ``Moment maps and surjectivity in various geometries,'' Conjecture \(6.16\), followed by Section \(6.5\) and Theorem \(6.17\): \url{https://aimath.org/WWN/momentmaps/momentmaps.pdf}. Page-number artifacts \(11,12\) and the Young-diagram extraction are not new problem labels. No PDF was opened in this drafting session.
\item Conjecture attribution as supplied: Nekrasov--Shatasvili--Moore 1998 and Hausel--Szenes 2004. Arithmetic attributions: Carlos Simpson; Hausel--Rodriguez--Villegas 2003; Frobenius--Schur 1907; Deligne's mixed Hodge structure. Full titles and authoritative bibliographic details, including the rank-two universal-sector source, are absent.
\item Supplied research attempt \(1\), \texttt{partial\_result}, run status \texttt{partially\_solved}, rigor label ``proved,'' novelty confidence ``low.'' Registry classification \texttt{PARTIAL-PROGRESS}, suggested difficulty \texttt{L4}, updated 14 August 2026. Original remarks and literature fields are empty; no proposer or year is supplied.
\end{itemize}
\end{document}
